\documentclass[11pt]{article}
\usepackage[margin=2.1cm]{geometry}
\usepackage{booktabs}
\usepackage{amsmath}
\usepackage{amssymb}
\usepackage{enumitem}
\usepackage{array}
\usepackage{listings}
\usepackage{xcolor}
\usepackage[T1]{fontenc}
\usepackage{lmodern}

\setlength{\parindent}{0pt}
\setlength{\parskip}{3.6pt}
\newcommand{\code}[1]{\texttt{#1}}
\newcommand{\hr}{\vspace{2pt}\hrule\vspace{6pt}}
\renewcommand{\arraystretch}{1.12}

\definecolor{kw}{RGB}{0,90,160}
\definecolor{cmt}{RGB}{110,110,110}
\lstdefinestyle{vhdl}{
  language=VHDL, basicstyle=\ttfamily\footnotesize,
  keywordstyle=\color{kw}\bfseries, commentstyle=\color{cmt}\itshape,
  showstringspaces=false, columns=fullflexible, breaklines=true,
  frame=leftline, framesep=6pt, xleftmargin=8pt, aboveskip=6pt, belowskip=6pt
}
\lstset{style=vhdl}

\begin{document}

\begin{center}
{\Large\bfseries EE\,362 --- Pre-Lab 4: Solution}\\[2pt]
{\large Datapath Support Components}\\[4pt]
Fall 2026 \quad$\cdot$\quad 5 points \quad$\cdot$\quad TA marking guide
\end{center}
\hr

{\small Marking: Q1 = 2 points, Q2 = 1 point, Q3 = 2 points.}

\vspace{4pt}\hr

\section*{Question 1 --- Extension by hand (2 points)}

\subsection*{(a) The four results}

\begin{center}
\begin{tabular}{c|ccc|l|l}
\toprule
\# & \code{Immed16} & \code{Xsel} & \code{Xlsft} & \code{Xtnd32} & as a number \\
\midrule
0 & \code{0x8000} & 0 (zero) & 0 (none) & \code{0x00008000} & $+32768$ \\
1 & \code{0x8000} & 1 (sign) & 0 (none) & \code{0xFFFF8000} & $-32768$ \\
2 & \code{0x0002} & 1 (sign) & 1 ($\times 4$) & \code{0x00000008} & $2 \times 4 = 8$ \\
3 & \code{0xFFFF} & 1 (sign) & 1 ($\times 4$) & \code{0xFFFFFFFC} & $-1 \times 4 = -4$ \\
\bottomrule
\end{tabular}
\end{center}

\textbf{How each one is obtained:}
\begin{itemize}[leftmargin=1.6em,itemsep=1pt]
\item Row 0: pad the top 16 bits with \code{0}.
\item Row 1: bit 15 of \code{0x8000} is \code{1}, so the top 16 bits become all
\code{1}s.
\item Row 2: bit 15 of \code{0x0002} is \code{0}, so sign-extension pads zeros;
shifting left by two multiplies by four.
\item Row 3: \code{0xFFFF} sign-extends to \code{0xFFFFFFFF} ($-1$), and shifting
left by two gives \code{0xFFFFFFFC} ($-4$).
\end{itemize}

\textbf{The last column is the check worth teaching:} sign-extension followed by
$\times 4$ still multiplies the \emph{value} by four, negatives included. If a
student's row 3 is not $-4$, their shift dropped or duplicated a bit.

\textbf{Common errors:} row 1 written as \code{0x0000FFFF} (they extended the
wrong end); row 3 written as \code{0x3FFFFFFF} (they shifted \emph{right}); rows
2--3 shifted by one instead of two.

\subsection*{(b) Same input, two different answers}

\begin{quote}
\textbf{Model answer:} Nothing in the 16 bits themselves decides it. The
\textbf{instruction} decides --- the opcode says whether that immediate is a
signed number or an unsigned bit pattern, and the decoder sets \code{Xsel}
accordingly.
\end{quote}

\textbf{This is the point of the whole question.} The bits carry no meaning on
their own; meaning is supplied by the instruction that uses them. In Lab 10 the
student writes the decoder that sets this bit, so the answer given here is the one
they will need again.

\subsection*{(c) Why \code{ADDI} sign-extends and \code{LUI} zero-extends}

\begin{quote}
\textbf{Model answer:} \code{ADDI}'s immediate is a \textbf{signed number} that
gets added to a register, so its value must survive widening ---
\code{ADDI r1,r1,-1} has to actually subtract one, which requires the top bits to
become \code{1}s. \code{LUI}'s immediate is a \textbf{bit pattern} to be placed in
the register, not a magnitude, so there is no sign to preserve and the remaining
bits are filled with zeros.
\end{quote}

Full credit needs the contrast \emph{number} versus \emph{bit pattern}. Accept a
student who says ``\code{ADDI} must handle negative offsets''.

\textbf{Consequence worth saying aloud:} get this backwards and the CPU's loop
counter counts the wrong way in Lab 12 --- and the ALU, the register file and the
memory will all be innocent.

\vspace{4pt}\hr

\section*{Question 2 --- Interface and resolution (1 point)}

\subsection*{The entity}

\begin{lstlisting}
entity ImmedXtnd is
    port (
        Immed16 : in  std_logic_vector(15 downto 0);
        Xsel    : in  std_logic;                       -- 0 = zero, 1 = sign
        Xlsft   : in  std_logic;                       -- 1 = shift left by 2
        Xtnd32  : out std_logic_vector(31 downto 0)
    );
end entity ImmedXtnd;
\end{lstlisting}

Check the widths: input 16, output 32, and \textbf{two single-bit controls}.

\subsection*{The resolution table}

\begin{center}
\begin{tabular}{c||ccc}
\toprule
& \code{0} & \code{1} & \code{Z} \\
\midrule
\code{0} & \code{0} & \textbf{\code{X}} & \code{0} \\
\code{1} & \textbf{\code{X}} & \code{1} & \code{1} \\
\code{Z} & \code{0} & \code{1} & \code{Z} \\
\bottomrule
\end{tabular}
\end{center}

Read it as three rules:
\begin{itemize}[leftmargin=1.6em,itemsep=1pt]
\item two drivers \textbf{agreeing} give that value;
\item one driving and one at \code{'Z'} gives the driven value --- \textbf{this is
how a bus works};
\item two drivers \textbf{disagreeing} give \code{'X'}.
\end{itemize}

\begin{quote}
\textbf{Therefore \code{'X'} in a waveform never means the simulator is confused.}
It means two things drove one wire and disagreed. From Lab 9 onward that is the
fastest diagnosis available.
\end{quote}

\subsection*{The one sentence: why not a Boolean equation?}

\begin{quote}
\textbf{Model answer:} Because the output range is $\{0, 1, Z\}$ and Boolean
algebra only has $\{0, 1\}$. \code{'Z'} is not a value --- it is the
\emph{absence} of a driver --- so there is no AND/OR/NOT expression that produces
it.
\end{quote}

Accept also: ``the value on the wire depends on the \emph{other} drivers, so it is
not a function of this component's inputs alone.'' That is a deeper answer and
deserves a comment in the margin.

\vspace{4pt}\hr

\section*{Question 3 --- Test plan (2 points)}

\subsection*{Model vector list}

\begin{center}
\small
\begin{tabular}{c>{\raggedright\arraybackslash}p{5.6cm}>{\raggedright\arraybackslash}p{7.0cm}}
\toprule
\# & Inputs & What it proves \\
\midrule
1 & \code{0x8000}, \code{Xsel}=0, \code{Xlsft}=0 & zero-extension of a value whose top bit is 1 \\
2 & \code{0x8000}, \code{Xsel}=1, \code{Xlsft}=0 & \textbf{sign-extension}: same input, different answer \\
3 & \code{0x7FFF}, \code{Xsel}=1 & sign-extension of a \emph{positive} value pads zeros, not ones \\
4 & \code{0xFFFF}, \code{Xsel}=1, \code{Xlsft}=1 & extend \textbf{then} shift, in that order: \code{0xFFFFFFFC} \\
5 & bus: one buffer \code{Sel}=1, the other \code{Sel}=0 & the enabled driver owns the wire; the other is \code{'Z'} \\
6 & bus: both \code{Sel}=0 & the wire floats: all \code{'Z'}, driven by nobody \\
7 & bus: both \code{Sel}=1 with \emph{different} data & \textbf{contention}: the wire shows \code{'X'} \\
\bottomrule
\end{tabular}
\end{center}

Six vectors earn full credit. Vectors 2 and 7 are the ones that must appear.

\subsection*{The order-of-operations trap}

Vector 4 is worth a sentence at the bench. \textbf{Extension happens first, then
the shift.} Shifting first and extending afterwards gives a different answer, and
both orders look reasonable in a sentence of English. The specification's table
settles it, which is why the table is written before the code.

\subsection*{Contention is resolved \emph{per bit}}

Worth telling strong students: if two drivers put \code{0xAAAAAAAA} and
\code{0xBBBBBBBB} on the same bus, the result is \textbf{not} all \code{'X'}. The
patterns agree on some bits and disagree on others, so the wire shows \code{'X'}
only where they differ. A test that checks for ``all X'' would fail on a correct
simulator. Check the bits that actually conflict.

\subsection*{The closing question: \code{0x0001} extends the same either way}

\begin{quote}
\textbf{Model answer:} A test value must be able to \emph{distinguish} the cases.
\code{0x0001} has a top bit of \code{0}, so sign- and zero-extension give the same
answer and the vector cannot fail --- it proves nothing about \code{Xsel}. Choose
values where the two behaviours differ, which here means \textbf{bit 15 set}.
\end{quote}

\textbf{This generalises, and it is the real lesson of the pre-lab:} a vector that
cannot fail is not a test. The same reasoning is why Lab 3's data patterns must
differ between inputs, and why Lab 2's toggle vector must move \code{D} between
edges.

\vspace{6pt}
\hr
{\small \textbf{At the bench:} any student whose Q1(a) row 1 is \code{0x0000FFFF}
or whose row 3 shifted right will build an \code{ImmedXtnd} that passes their own
testbench and breaks Lab 12. Correct it before they start.}

\end{document}
